\section{Experimental setup}
\label{sec:setup}

% ROUGH OUTLINE -- one bullet = one point (roughly one paragraph in the final
% text). All concrete numbers (grids, priors, sensor counts, swept values) go
% into the tables during writing.

\subsection{Urban test cases}
\label{sec:setup:cases}
\begin{itemize}
  \item Two cases bracketing the idealized--realistic axis
        (\tabref{tab:cases}):
        \begin{itemize}
          \item \textbf{Case A:} Xie \& Castro staggered cube array ---
                canonical, well-documented benchmark \cite{xie2006les,xie2009modelling}.
          \item \textbf{Case B:} ``Barcelona'' real urban block --- irregular
                heights and street layout from STL geometry.
        \end{itemize}
  \item Boundary conditions (both cases): prescribed inflow profile
        (\velmag, \inflowangle, power-law shear \shearexp), periodic lateral,
        free-slip lid, no-slip walls; spin-up before observations are drawn.
  \item Grids and near-wall treatment deliberately differ per solver
        (\tabref{tab:solver-differences}) --- part of the intended \discrepancy.
\end{itemize}

\begin{table}[t]
  \centering
  \caption{Test cases: domain, grid, and boundary conditions per solver.}
  \label{tab:cases}
  \todo{table body}
\end{table}

\begin{table}[t]
  \centering
  \caption{Structural solver differences (discretization, closure, geometry
    representation, wall treatment) --- the source of \discrepancy.}
  \label{tab:solver-differences}
  \todo{table body}
\end{table}

\subsection{Cross-model experimental design}
\label{sec:setup:crossmodel}
%==============================================================================
%  figures/cross_model_setup.tex
%  The cross-model experimental pipeline: a truth model generates synthetic
%  observations that a DIFFERENT assimilation model never sees during ES-MDA.
%  Adapted from the "Setup" slide of presentation.tex.
%  \input where the figure should float (experimental setup / cross-model).
%==============================================================================
\begin{figure}[tb]
  \centering
  \begin{tikzpicture}[font=\scriptsize,>=latex,scale=1.0,transform shape,
      every node/.append style={align=center},
      obox/.style={rounded corners=3pt,draw=uaOrange,line width=0.9pt,fill=uaOrange!6,
                   inner sep=4pt,text width=1.9cm,minimum height=0.85cm},
      tbox/.style={rounded corners=3pt,draw=uaTeal,line width=0.9pt,fill=uaTeal!6,
                   inner sep=4pt,text width=1.9cm,minimum height=0.85cm},
      ar/.style={->,line width=0.9pt}]
    % truth lane (top)
    \node[obox] (t1) at (0,1.5)    {truth inflow\\$\param(t)$};
    \node[obox] (t2) at (2.45,1.5) {truth\\solver};
    \node[obox] (t3) at (4.90,1.5) {truth\\state};
    \node[obox] (t4) at (7.35,1.5) {obs.\\operator};
    \node[obox,fill=uaOrange!16] (t5) at (9.80,1.5) {obs.\ $\obs$\\{\scriptsize($+$ noise)}};
    \draw[ar,uaOrange] (t1)--(t2);
    \draw[ar,uaOrange] (t2)--(t3);
    \draw[ar,uaOrange] (t3)--(t4);
    \draw[ar,uaOrange] (t4)--(t5);
    % assimilation lane (bottom)
    \node[tbox] (a1) at (0,-1.5)
      {prior inflow\\ensemble\\{\scriptsize($\Ne$)}};
    \node[tbox,fill=uaTeal!10,text width=6.4cm,minimum height=1.1cm] (a2) at (5.2,-1.5)
      {\textbf{\color{uaTeal}\esmda}\\[2pt]
       {\scriptsize forecast $\to$ predict obs.\ $\to$ update $\to$ repeat}};
    \node[tbox] (a3) at (9.80,-1.5) {posterior\\inflow};
    \draw[ar,uaTeal] (a1)--(a2);
    \draw[ar,uaTeal] (a2)--(a3);
    % observations feed INTO ES-MDA
    \draw[ar,uaCharcoal] (t5.south) --
       node[pos=0.62,right=1pt,font=\scriptsize]{observations} (a2.north east);
    % rollout loop
    \draw[ar,uaTeal] (a3.south) .. controls (9.8,-2.75) and (5.2,-2.75) ..
       node[pos=0.5,below,font=\scriptsize]{rollout: warm-start state, extrapolate prior} (a2.south);
  \end{tikzpicture}
  \caption{Cross-model experimental design. A \truthmodel\ (orange, top)
    generates synthetic observations $\obs$ that are corrupted with noise and
    fed to \esmda; the \assimmodel\ (teal, bottom) is a \emph{different} solver
    and never sees the data-generating model, so its structural \discrepancy\ is
    real and irreducible. This is the setup that avoids the \invcrime\ of
    same-model \twinexp{s}. \todo{confirm ensemble size and window count shown.}}
  \label{fig:cross_model_setup}
\end{figure}
   % \figref{fig:cross_model_setup}
\begin{itemize}
  \item Protocol (\figref{fig:cross_model_setup}): \truthmodel\ generates
        synthetic observations with fixed true parameters; a \emph{different}
        \assimmodel\ runs \esmda; recovery is scored against the known truth.
  \item Pairings (\tabref{tab:pairings}): diagonal $=$ same-model \twinexp\
        baselines (\invcrime); off-diagonal $=$ \crossmodel, ordered from
        ``near'' (LES$\to$LES) to ``far'' (LES$\leftrightarrow$\lbmmodel)
        \discrepancy.
  \item True parameters drawn independently of the assimilation prior (prior
        not secretly centred on truth).
\end{itemize}

\begin{table}[t]
  \centering
  \caption{Truth $\times$ assimilation pairings ($3\times3$; diagonal $=$
    \twinexp\ baseline, off-diagonal $=$ \crossmodel).}
  \label{tab:pairings}
  \todo{table body}
\end{table}

\subsection{Estimated parameters and priors}
\label{sec:setup:params}
\begin{itemize}
  \item Estimated vector: forcing (\inflowangle, \velmag, \shearexp) $+$
        closure knob \sgs\ ($+$ \pgrad\ for \pyudales\ only); Gaussian priors
        with physical bounds, static or AR(2) time-varying (\tabref{tab:priors}).
  \item Key design point: \shearexp\ and \sgs\ are \discrepancy-absorbing
        knobs --- nearly redundant in a \twinexp, load-bearing cross-model;
        their admissible ranges are retuned per assimilation solver for
        numerical stability.
\end{itemize}

\begin{table}[t]
  \centering
  \caption{Parameters: priors, bounds, truth values, temporal model.}
  \label{tab:priors}
  \todo{table body}
\end{table}

\subsection{Observation network}
\label{sec:setup:obs}
\begin{itemize}
  \item Street-level velocity sensors, split into an assimilation set and a
        held-out validation set (\figref{fig:sensors}); observations
        time-aggregated over intervals; noise level $\sigma$ set to avoid the
        unobservable-observation regime (and swept in the sensitivity study).
\end{itemize}

\begin{figure}[t]
  \centering
  \todo{plan-view sensor map per case (assimilation vs.\ held-out)}
  \caption{Sensor network for the two cases.}
  \label{fig:sensors}
\end{figure}

\subsection{Sensitivity study design}
\label{sec:setup:study}
\begin{itemize}
  \item One-factor-at-a-time sweep from a fixed cross-model baseline
        (\tabref{tab:study}). Core factors: model pairing, observation error
        $\sigma$, estimated-parameter set (with/without \shearexp, \sgs),
        static vs.\ AR(2). Secondary: \Ne, \Na, \Nw, localization, inflation,
        sensor count/placement, prior spread, window length.
\end{itemize}

\begin{table}[t]
  \centering
  \caption{Sensitivity factors: swept values, baseline, hypothesis tested.}
  \label{tab:study}
  \todo{table body}
\end{table}

\subsection{Evaluation protocol}
\label{sec:setup:eval}
\begin{itemize}
  \item Per configuration: parameter \RMSE/\CRPS\ (prior vs.\ posterior),
        calibration (spread, coverage), and predictive skill at held-out
        sensors; repeated over seeds; \twinexp\ as ceiling, prior as floor;
        forcing and closure parameters reported separately.
\end{itemize}
